\newpage
\section{Estimación del tiempo hasta velocidad vertical cero}
Una vez obtenida $V_{r,k}$, podemos estimar cuánto tiempo falta para que la velocidad vertical llegue a cero. Durante el ascenso, este instante corresponde al apogeo si después comienza el descenso. Que la velocidad vertical sea cero no significa que la velocidad lateral también lo sea.

Usaremos las ecuaciones del movimiento rectilíneo uniformemente acelerado (MRUA), aplicadas a la componente vertical. Este modelo supone una aceleración vertical constante durante el intervalo que se quiere predecir. El movimiento completo del cohete puede tener componentes laterales.

\subsection{Tiempo restante a partir de la velocidad actual}
Sea $\Delta t_0$ el tiempo restante desde el instante actual $t_k$, y $a_{\mathrm{pred},k}$ la aceleración vertical que suponemos constante para la predicción. La ecuación de velocidad es
\begin{equation}
 V_z(t_k+\Delta t)=V_{r,k}+a_{\mathrm{pred},k}\Delta t.
\end{equation}
Al imponer $V_z=0$ y despejar el tiempo,
\begin{equation}
 0=V_{r,k}+a_{\mathrm{pred},k}\Delta t_0,
 \qquad
 \boxed{\Delta t_0=-\frac{V_{r,k}}{a_{\mathrm{pred},k}}}.
\end{equation}
Durante un ascenso con desaceleración, $V_{r,k}>0$ y $a_{\mathrm{pred},k}<0$, por lo que el tiempo calculado es positivo. El instante estimado de apogeo es $t_{\mathrm{ap}}=t_k+\Delta t_0$; no debe confundirse con el tiempo restante.

Si $a_{\mathrm{pred},k}=0$, el modelo no predice una detención futura para una velocidad inicial positiva. Si $a_{\mathrm{pred},k}>0$, el resultado algebraico es negativo y no representa un apogeo futuro: bajo esa hipótesis el cohete continúa aumentando su velocidad vertical. Una aceleración cercana a cero hace que la predicción sea muy sensible al ruido.

\subsection{Aproximación después de terminar el empuje}
Si el empuje ya terminó y se desprecia la resistencia del aire, la aceleración vertical es $a_{\mathrm{pred},k}=-g$. Entonces,
\begin{equation}
 \boxed{\Delta t_0=\frac{V_{r,k}}{g}}.
\end{equation}
Por ejemplo, si en un instante de vuelo sin empuje se estima la misma velocidad del ejemplo anterior, $V_{r,k}=8.328\,\mathrm{m/s}$, se obtiene
\begin{equation}
 \Delta t_0=\frac{8.328}{9.81}\approx0.849\,\mathrm{s}.
\end{equation}
Este cálculo corresponde a una hipótesis nueva de vuelo sin empuje. En las lecturas del ejemplo anterior, $a_{z,2}=+2.662\,\mathrm{m/s^2}$: el cohete todavía aumenta su velocidad vertical. Mantener esa aceleración positiva no permite predecir un apogeo futuro, y no sería correcto cambiar su signo para obtener un tiempo positivo.

En vuelo real, el arrastre y cualquier empuje residual modifican la aceleración. Como aproximación local puede emplearse una aceleración vertical medida y filtrada, siempre que sea negativa y razonablemente estable; no se garantiza que permanezca constante hasta el apogeo. El tiempo se recalculará con cada nueva estimación válida de velocidad y aceleración. Llegar a una predicción de tiempo cero no sustituye la confirmación del apogeo mediante las mediciones.
\end{document}
